\section{Introduction}
A renewal provider has accepted a total obligation but must implement it through a limited catalog of terminal commands. Different histories have different expected caps, realization ceilings, and costs of intermediate service. The provider chooses both the installed commands and the allocation of the obligation across histories. Opening charges discourage unnecessary commands. A useful algorithm must optimize these decisions jointly: changing the obligation or adding commands to repair an approximation changes the problem the provider agreed to solve.

This paper identifies two exact path representations of that joint problem. The first is a primal selected-boundary representation for a common increasing concave terminal reward. Every neighboring pair of selected commands owns a terminal-reward layer and the service costs of the histories whose eligibility ends at that boundary. These layers admit an explicit history-level reconstruction. Their capacities sum to the original available obligation, regardless of how the command book partitions histories. A uniform resource approximation can consequently be repaired at the original promise and command budget.

A history-level component-packing argument extends the primal representation to heterogeneous rewards by replacing a common chord with a concave within-edge allocation profile. The second representation is an exact priced path. It permits different increasing concave terminal rewards across histories. At a fixed obligation price, the optimal response of each history separates into its positive chord increments and its terminal service support. These terms can again be assigned to selected boundaries. One ordered dynamic program then optimizes all feasible anchors and command books. It computes a genuine original-space Lagrangian bound, not a bound on a gridded or relaxed-feasibility surrogate. The priced identity does not require different histories to share an ordering of marginal rewards; each history only needs its own concave ordering.

We first establish a complexity frontier. The exact design decision problem is NP-complete even with common unit linear reward, zero service costs, uniform history weights and a nonbinding command budget. A direct subset-sum construction gives the proof; the obstruction is not inferred from the running time of a particular algorithm. In contrast, common expected caps admit an exact two-command reduction under a common terminal reward, even when the full catalog has many eligibility classes. Enumerating singletons and pairs solves this subclass in polynomial time without an accuracy parameter. It is distinct from the retained exact common-prefix subclass: neither identical caps nor identical eligibility is imposed by the general model. For general caps, rational price-path bounds support a finite, grid-free branch-and-bound algorithm. Every explored or unexplored leaf carries a valid upper bound, and every incumbent is an original-space lottery policy. The resulting interval is valid even when a run reaches its resource limit. The unrestricted exact search can still be exponential. The improvement is a polynomial node oracle and compact, data-dependent certificates, not an assertion of polynomial exact optimization in the NP-complete regime.

The path representation also strengthens preprocessing. A price path forced to contain a given command bounds every book using that command, including books in which it is the anchor. Comparing this bound with a feasible incumbent gives safe exclusion. This is complementary to the retained, inexpensive non-anchor envelope, which preserves targets uniformly but ignores the accepted promise and budget. Neither test guarantees a positive removal rate. We therefore measure exclusion rates, threshold slacks, missed exclusions, anchor removals, and total computation rather than identify an unchanged incumbent as proof that a new command is irrelevant.

The closest algorithmic antecedents are Lagrangian constrained-path methods, not a new form of generic branching. \citet{HandlerZang1980} combine dual bounds with path enumeration, and \citet{BeasleyChristofides1989} use Lagrangian bounds for tree search and problem reduction. \citet{Fisher1981} explains the general integer-programming role of such relaxations. Those mechanisms are used here, not claimed as inventions. The necessary additional theorem is that the uneliminated renewal problem has the stated exact priced arc functions, despite individual caps, pre-draw service, realization ceilings, and endogenous command choices. The proof also supplies a heterogeneous-reward extension that the common-reward primal rearrangement cannot provide.

The retained resource approximation relates to constrained-path discretization \citep{Hassin1992,LorenzRaz2001}, capped service allocation to separable convex resource allocation \citep{Patriksson2008}, and the uniform-grid acceleration to totally monotone matrix search \citep{Aggarwal1987}. Ordered quantizer design \citep{Wu1991} and adjacent mean-preserving randomization \citep{PagesWilbertz2012} supply related ordering and interpolation ideas. They do not by themselves allocate a fixed aggregate obligation jointly with a paid, limited command book and history-specific realization restrictions. Our novelty claims concern the two exact decompositions, their constructive connection to the original constraints, and the hardness and common-cap tractability results. Gridding, dual weak bounds, and incumbent-based exclusion retain their established interpretations.

The computational study separates theorem validation from solver evaluation. Exact regressions compare against every feasible book on small models and reject altered certificates. A matched-limit study varies histories, catalog size, budget, accuracy, charges, and reward heterogeneity. It includes cases outside the allotted enumeration range, records interrupted runs, and reports proof-verification cost separately. Historical fine-accuracy cases are rerun rather than compared to timings from another machine. Additional zero-width runs expose anytime behavior, and near-threshold screening cases compare the two safe rules with exact inclusion relevance. Every operational input remains synthetic. The paper makes a methodological case; it does not infer field calibration or a universal speed advantage from those experiments.
